% !PRINTED 465
\weberpage{465}

%% [running head, not body text -- commented out 2026-08-07] Page furniture (running
%% head, folio, signature) is excluded from this edition by standing policy.
%% \noindent \S\,126. \hfill The class invariants $f(\omega)^{24}$. \hfill 465

%<<CONT
1. If $D\equiv 4\pmod{8}$, we set
\[
n=m,\quad x=1,\quad y=0,
\]
\[
\alpha=\frac{B}{2},\quad \beta=A,
\]
\[
\gamma\equiv -C,\quad \delta\equiv\frac{B}{2}\pmod{2}.
\]

The congruence (9) therefore requires, since not all three coefficients $A$, $B$, $C$ can be even, that the two outer coefficients be odd and the middle one divisible by 4; and this occurs only for one of the three forms (3), (5). Consequently one of the three functions (2) satisfies equation (7), and is therefore a class invariant.

\[
\begin{aligned}
2.\quad &D\equiv 0\pmod{8},\\
&n=m+1,\quad x=1,\quad y=\pm 2,\\
&\alpha=\frac{B}{2}-1,\quad \beta=A,\\
&\gamma\equiv -C,\quad \delta\equiv\frac{B}{2}+1\pmod{2}.
\end{aligned}
\]

Here too, then, the two outer coefficients must be odd and $B\equiv 2\pmod{4}$, and everything then stands as above. With this it is proved:

1. \gesperrt{If $\omega$ is the root of a primitive quadratic equation with negative discriminant $D$ divisible by 4 and with odd outer coefficients, then $f(\omega)^{24}$ is a class invariant for the discriminant $D$.}

The assumption that $A$ is relatively prime to $n$, i.\,e. to $m+1$ or to $m$, may afterwards be given up as unessential.

For if we apply a linear transformation $(S)$ to $\omega$, then, for odd $A$, $C$, equation (3) passes into an equivalent equation in which the outer coefficients are both odd then and only then, when $(S)$ belongs to the first or the second class, that is, when $f(\omega)^{24}$ remains unchanged under $(S)$. One can then further determine $(S)$ so that the first coefficient in the transformed equation (3) becomes relatively prime to any number one pleases, and hence also to $n$.
